\subsection{Inverse of a Power Function} 
Let's look for a pattern in the inverses of power functions.
\begin{itemize}
\item If $n$ is a whole number, the inverse of $f(x)=x^n$ is $f^{-1}(x)=\makeblank[1in]$. If $n$ is \makeblank[1in], we need to \makeblank to \makeblank[1in].
\item What is the inverse of $f(x)=x^{-1}$?
\end{itemize}
\begin{lpic}{}{4}
\regtextleft{1}{-4}{$f^{-1}(x)=$};
\end{lpic}
\begin{itemize}
\item If $\frac{m}{n}$ is a positive fraction in lowest terms, what is the inverse of $f(x)=x^{m/n}$?
\end{itemize}
\begin{lpic}{}{5}
\regtextleft{1}{-5}{$f^{-1}(x)=$};
\end{lpic}
\begin{itemize}
\item Note:
\begin{itemize}
\item If $n$ is \makeblank[1in], the domain of $f$ is \makeblank[1.5in]
\item If $m$ is \makeblank[1in], then $f$ is not\makeblank[1.5in], so we restrict the domain of $f$ to \makeblank[1.5in] to have an inverse.
\end{itemize}
\end{itemize}
Looking at the pattern here, we can make a general statement for all power functions:
\[
\mbox{If }
f(x)=x^p
\mbox{ then }
f^{-1}(x)=\makeblank[1in]
\]
We can use this fact to solve equations.